\chapter{Multilevel Moderated Networks}
Experience sampling approaches and access to intensive longitudinal data have been important components of the popularity of these types of models. This has opened up possibilities of applying time series methods and investigating both linear and nonlinear relationships among day-to-day experiences such as mood and vulnerability to depression. Experience sampling methods generally involve having participants complete surveys and complete psychological measures throughout their daily life---this may occur only once or twice per day (as in the traditional `daily-diary' study), or anywhere from 5 to 10 times per day [CITE]. These time-intensive data are available in much greater volume now given the widespread accessibility of mobile phones and other `smart' devices that pervade our daily lives. These methods appear frequently in the psychological networks literature, and provide a potentially powerful view into intraindividual dynamics as well as those that can be observed between-subjects. 

\section{Within-Person Processes}
Only longitudinal data allow for the proper separation of between-person and within-person effects, and this is critically needed for evaluating many theories in psychology \cite{curran2011disaggregation}. However, it is sometimes difficult to fully articulate the ways in which a given influence on an outcome might vary in magnitude and form when looking within versus across persons. Indeed, the influences at these levels may be operating simultaneously---and even in opposite directions. Although repeated measures data are becoming more common, more emphasis is needed on methods for better testing our theories and hypotheses \cite{curran2011disaggregation}. 

Data collected from multiple individuals at a single time point can only be used for thinking about between-person processes. Repeated measures taken from multiple individuals, however, contain information about both levels of interest. Still, we must employ methods carefully in order to prevent confounding the two sources of variability. Take the following example: research has shown that an individual is more likely to experience a heart attack while exercising, but at the same time people who exercise tend to have a lower risk of heart attack \cite{curfman1993exercise,mittleman1993triggering}. Generalizing either effect to the other source would lead to an error of inference. Moreover, only studying one of the two sources would limit a more complete understanding of the relationship between exercise and heart attack risk. The ecological fallacy is an example of reasoning incorrectly by generalizing between-person effects to the level of the individual; even if the relations at each level are the same, the relation at one level is neither necessary nor sufficient to imply the same relation at the other level. 

\subsection{Multilevel Models}
We denote the repeated measure observed at time point $t$ for individual $i$ as $y_{ti}$. For the linear growth model the observed repeated measure is expressed as a simple linear function of time: $y_{ti} = \beta_{0i} + \beta_{1i} x_{ti} + r_{ti}$, where the coefficients are the intercept and linear slope for individual $i$, $x_{ti}$ is the observed value of time at assessment $t$ for individual $i$, and $r_{ti}$ is the time- and individual-specific residual. This represents the within-person trajectory, and is also referred to as the level-1 equation. 

In growth models, values of the intercept and slope components vary randomly across persons: $\beta_{0i} = \gamma_{00} + u_{0i}$, and $\beta_{1i} = \gamma_{10} + u_{1i}$. Here, $\gamma_{00}$ and $\gamma_{10}$ are the overall mean intercept and slope, and $u_{0i}$ and $u_{1i}$ are the person-specific deviations from those means. This captures the between-person differences in within-person change and is referred to as the level-2 equation. The formal statistical model then results from the substitution of equation 2 into equation 1 that in turn defines the reduced-form expression
\begin{equation}
    y_{ti} = (\gamma_{00} + \gamma_{10} x_{ti}) + (u_{0i} + u_{1i} x_{ti} + r_{ti}).
\end{equation}
Time-invariant covariates $w_i$ would be included at level 2: $\beta_{0i} = \gamma_{00} + \gamma_{01} w_i + u_{0i}$, and $\beta_{1i} = \gamma_{10} + \gamma_{11} w_i + u_{1i}$. Time-varying covariates, however, would be included at level 1, as they vary over both individual and time point. We denote the TVC as $z_{ti}$. TVCs simultaneously contain both within-person and between-person variability, for example: $z_{ti} = \bar{z}_i + r_{ti}$, where $\bar{z}_i$ is the person-specific mean of the TVC, pooling over time, and $r_{ti}$ is the time-specific deviation of the TVC from the person-specific mean. 

Take this example: we want to know if daily fluctuations in anxiety $z$ predict substance use $y$.
\begin{equation}
    \begin{split}
        y_{ti} &= \beta_{0i} + \beta_{1i} z_{ti} + r_{ti} \\
        \beta_{0i} &= \gamma_{00} + u_{0i}, \quad \beta_{1i} = \gamma_{10} \\
        y_{ti} &= (\gamma_{00} + \gamma_{10} z_{ti}) + (u_{0i} + r_{ti})
    \end{split}
\end{equation}
Here we have a random intercept, but no random slope. Even so, the fixed effect of the TVC on the outcome $\gamma_{10}$ represents an aggregation of within-person and between-person influences of the TVC on the outcome. The reason is that $z_{ti}$ varies between individuals in mean level, and within individuals across time. Thus, when we estimate just one effect of $z_{ti}$, the result is a combination of potentially different effects operating at two different levels of analysis. So, we must decompose the TVC into components that isolate between- and within-person differences in order to differentiate their effects. 

\subsubsection{Sources of variability in two-level models}
To quickly recap: in two-level models we partition the outcome variable into two orthogonal sources of variability---level-2 variability refers to cluster mean differences on the outcome variable, and level-1 variability refers to within-cluster differences on the outcome variable. Consider the example of a study measuring chronic pain in which daily observations (level 1) are nested within individuals (level 2). Suppose the researchers are interested in predicting participants' daily affect ratings. Level-2 variability refers to the person-to-person differences in average affect levels (e.g., some people have higher levels of positive affect than others), while level-1 variability refers to day-to-day fluctuations around individual participants' average affect levels.

The unconditional model with no predictors is: $Y_{ij} = \gamma_{00} + u_{0j} + \varepsilon_{ij}$, where $\gamma_{00}$ is the weighted grand mean, $u_{0j}$ is the residual that represents cluster mean differences on the outcome variable, and $\varepsilon_{ij}$ represents differences between scores and their cluster-specific means.$^3$ I.e., $u_{0j}$ is the difference between participant $j$'s average affect level and the weighted grand mean, and $\varepsilon_{ij}$ is the difference between participant $j$'s affect rating on day $i$ and her average affect level. If we wanted to predict daily affect from daily sleep ratings, we would have a two-level model: $Y_{ij} = \gamma_{00} + \gamma_{10} X_{ij} + u_{0j} + \varepsilon_{ij}$.

\emph{In Gina Mazza's notation, we use $\varepsilon_{ij}$ and $\sigma_{\varepsilon}^2$ (rather than $r_{ij}$ and $\sigma^2$) to represent the level-1 residual and its variance, as well as $\sigma_{u_{0j}}^2$, $\sigma_{u_{1j}}^2$, etc. (rather than $\tau_{00}$, $\tau_{11}$, etc.) to represent the level-2 residual variances.}

In two-level models, all level-1 predictors potentially have two sources of variability, and thus can have within-cluster and/or between-cluster associations with the outcome variable. In the previous example, for instance, a within-cluster association between daily sleep and daily affect means that day-to-day fluctuations around a participants' average sleep level predict day-to-day fluctuations around his/her average affect level; whereas a between-cluster association means that a participants' average sleep level predicts his/her average affect level. Because here we are representing both the within- and between-cluster associations with one level-1 regression coefficient $\gamma_{10}$, we assume that the within- and between-cluster associations between the level-1 predictor and the outcome are equal. In fact, the level-1 regression coefficient $\gamma_{10}$ is actually a weighted average of the within- and between-cluster associations between $X_{ij}$ and $Y_{ij}$. 

If $b_T$ stands in for any level-1 regression coefficient: $b_T = \eta_X^2 b_B + (1 - \eta_X^2)b_W$, where $b_B$ and $b_W$ are the between- and within-cluster effects of the level-1 predictor on the outcome, and $\eta_X^2$ is the ratio of the between-cluster sum of squares on the level-1 predictor to the total sum of squares on the level-1 predictor. Based on this equation, the level-1 regression coefficient $\gamma_{10}$ only unambiguously estimates a level-specific association if: (a) the within- and between-cluster associations are equal, (b) there is no variability at level 2 (i.e., $\eta_X^2 = 0$), or (c) there is no variability at level 1 (i.e., $\eta_X^2 = 1$). Level-2 predictors, however, only have level-2 variability. 

\subsection{Person-Mean Centering}
These effects can be unambiguously disaggregated within the multilevel model using person-mean centering. However, with time-varying covariates repeated measures are nested within each individual, and so there are two means to consider: the grand mean of the TVC pooling over all time points and all individuals, and each person-specific mean pooling over all time points within the individual. 
\begin{equation}
    \ddot{z}_{ti} = z_{ti} - \bar{z}_{..}
\end{equation}
where $\ddot{z}_{ti}$ is the grand mean centered TVC, $z_{ti}$ is the observed TVC, and $\bar{z}_{..}$ is the grand mean of $z_{ti}$ pooling over all individuals and all time points. Then, we can calculate the person-specific mean of the TVC for each individual
\begin{equation}
    \dot{z}_{ti} = z_{ti} - \bar{z}_i
\end{equation}
where $\dot{z}_{ti}$ is the person-mean centered TVC, and $\bar{z}_i$ is the person-specific mean. Now, we can use $z_{ti}$, $\dot{z}_{ti}$, or $\ddot{z}_{ti}$ as the level-1 predictor, and each is associated with a potentially different inference with respect to the disaggregation of effects. 

Direct estimates of the between-person and within-person effects can be obtained by incorporating the person-mean centered TVC at level-1 ($\dot{z}_{ti}$), and the person-mean at level-2 ($\bar{z}_i$). 
\begin{equation}
    \begin{split}
        y_{ti} &= \beta_{0i} + \beta_{1i} \dot{z}_{ti} + r_{ti} \\
        \beta_{0i} &= \gamma_{00} + \gamma_{01} \bar{z}_i + u_{0i} \\
        \beta_{1i} &= \gamma_{10}
    \end{split}
\end{equation}
The reduced form equation for this model is:
\begin{equation}
    y_{ti} = (\gamma_{00} + \gamma_{01} \bar{z}_i + \gamma_{10} \dot{z}_{ti}) + (u_{0i} + r_{ti})
\end{equation}
where $\gamma_{00}$ is the intercept or grand-mean, $\gamma_{01}$ is a direct estimate of the between-person effect, and $\gamma_{10}$ is a direct estimate of the within-person effect. It's also important to note that person-mean centering can lead to an underestimation of the autoregressive effects \cite{hamaker2015center,bringmann2016assessing}.

\subsubsection{Other types of centering}
There are two centering options for level-2 predictors: raw score scaling (RAS), and grand mean centering (CGM). There are three centering options for level-1 predictors in two-level models: RAS, CGM, and centering within context (CWC; also called centering within clusters or group mean centering). This notation comes from \cite{kreft1995effect}, which is a seminal work on centering in multilevel models. 

RAS means simply leaving the predictor uncentered. CGM deviates scores around the grand mean: $X_{CGM} = X_{ij} - \bar{X}$. Because CGM deviates scores around the same constant, it preserves level-1 and level-2 variability in the level-1 predictor. CWC deviates scores around their cluster-specific means: $X_{CWC} = X_{ij} - \bar{X}_j$. This makes it so that all clusters have a mean of zero after centering; thus, there is no variability in the cluster means of the centered scores (i.e., no level-2 variability) after CWC. However, introducing the cluster means of the level-1 predictor as a level-2 predictor allows for a contextual effect: 
\begin{equation}
    Y_{ij} = \gamma_{00} + \gamma_{10} X_{ij} + \gamma_{01} \bar{X}_j + u_{0j} + \varepsilon_{ij}
\end{equation}
With CWC, the level-1 predictor $X_{ij}$ and cluster means $\bar{X}_j$ are uncorrelated. Thus, $\gamma_{10}$ quantifies the within-cluster association and $\gamma_{01}$ quantifies the between-cluster association. If we wanted to allow for the effect of $X_{ij}$ to differ across clusters, we could allow for a random slope: $Y_{ij} = \gamma_{00} + \gamma_{10} X_{ij} + \gamma_{01} \bar{X}_j + u_{0j} + u_{1j} X_{ij} + \varepsilon_{ij}$. 

\subsubsection{Level 1 interaction}
This is an interaction between level-1 predictors:
\begin{equation}
    Y_{ij} = \gamma_{00} + \gamma_{10} X_{ij} + \gamma_{20} Z_{ij} + \gamma_{30} X_{ij} Z_{ij} + u_{0j} + u_{1j} X_{ij} + u_{2j} Z_{ij} + u_{3j} X_{ij} Z_{ij} + \varepsilon_{ij}
\end{equation}
where $\gamma_{30}$ is the regression coefficient for the level-1 interaction, $u_{1j}$ is the residual that allows the effect of $X_{ij}$ to differ across clusters, $u_{2j}$ is the residual that allows the effect of $Z_{ij}$ to differ across clusters, and $u_{3j}$ is the residual that allows the effect of the interaction $X_{ij} Z_{ij}$ to differ across clusters.

With RAS or CGM, a level-1 interaction potentially yields a composite product term $X_{ij} Z_{ij}$ with four sources of variability. 

\subsection{Multilevel Vector Autoregression}
Multilevel VAR models assume stationarity of the modeled process, as trends in the data may lead to spurious correlations between lagged variables \cite{rovine2006multilevel}. At level one observations are modeled as
\begin{equation}
    \begin{split}
        y_{ti} &= \mu_i + a_{ti} \\
        \mu_i &= \mu + u_{0i} \sim \mathcal{N}(0, \sigma_{u0}^2) \\
        a_{ti} &= \phi_i a_{t-1,\,i} + e_{ti} \sim \mathcal{N}(0, \sigma_{e}^2) \\
        \phi_{i} &= \phi + u_{1i} \sim \mathcal{N}(0, \sigma_{u1}^2) \\
        y_{ti} &= (\mu + \phi a_{t-1,\,i}) + (u_{0i} + u_{1i} a_{t-1,\,i} + e_{ti})
    \end{split}
\end{equation}

\subsubsection{Making the model compatible with software}
Here we use an individually-centered lagged autoregressive parameter $(y_{t-1,\,i} - \mu_i)$ such that the model becomes
\begin{equation}
    \begin{split}
        y_{ti} &= \mu_i + \phi_i^c (y_{t-1,\,i} - \mu_i) + e_{ti} \\
        \mu_i &= \gamma_{00}^c + u_{0i}^c \\
        \phi_i^c &= \gamma_{10}^c + u_{1i}^c
    \end{split}
\end{equation}

\subsubsection{Simple example of random effects}
This is how we can think about random effects in a multilevel VAR model:
\begin{equation}
    \begin{split}
        Y_{ptj} &= \gamma_{0pj} + \gamma_{1pj} Y_{1,\,p,\,t-1} + \gamma_{2pj} Y_{2,\,p,\,t-1} + \gamma_{3pj} Y_{3,\,p,\,t-1} \\
        \gamma_{kpj} &= \beta_{kj} + b_{kpj}
    \end{split}
\end{equation}
where $\beta_{kj}$ represents the fixed effect of variable $k$ on variable $j$ over all individuals, and $b_{kpj}$ represents the person-specific deviation from this average effect.

\subsection{Multilevel Graphical VAR}
This allows us to study multivariate intraindividual dynamics as well as between-subjects overlap and differences. We assume that the number of time points $t \in \{1,2,\ldots,T\}$ may differ across individuals, and that measurement occasions are nested within individuals $i \in \{1,2,\ldots,N\}$. The level-1 equations are then:
\begin{equation}
    \begin{split}
        \matr{y}_{t,\,i} &= \bm{\mu}_i + \bm{\beta}_i (\matr{y}_{t-1,\,i} - \bm{\mu}_i) + \bm{\varepsilon}_{t,\,i} \\
        \bm{\varepsilon}_{t,\,i} &\sim \mathcal{N}(\matr{0}, \bm{\Theta}_i) \\
        \bm{\Theta}_i^{-1} &= \bm{K}_i^{\bm{\Theta}}
    \end{split}
\end{equation}
where $\bm{\mu}_i$ is the stationary mean vector of subject $i$ (we can no longer assume within-subjects means are zero without a loss of generality), $\bm{\beta}_i$ encodes the person-specific temporal network, and $\bm{K}_i^{\bm{\Theta}}$ encodes the person-specific contemporaneous network. 

For the level-2 equations, let $\bm{\beta}_{*}$ and $\bm{K}_{*}^{\bm{\Theta}}$ encode the expected temporal and contemporaneous networks when selecting a person at random---i.e., these are the between-subjects networks. We can assume without loss of generality that the data are grand-mean centered. 
\begin{equation}
    \begin{split}
        \mathbb{E} \lbrack \bm{\mu}_N \rbrack &= \matr{0} \\
        \mathbb{E} \lbrack \bm{\beta}_N \rbrack &= \bm{\beta}_{*} \\
        \mathbb{E} \lbrack \bm{K}_N^{\bm{\Theta}} \rbrack &= \bm{K}_{*}^{\bm{\Theta}}
    \end{split}
\end{equation}
Here, $\bm{\beta}_{*}$ and $\bm{K}_{*}^{\bm{\Theta}}$ encode the average parameters in the population, i.e., the fixed effects. Deviations from the fixed effects $(\bm{\beta}_i - \bm{\beta}_{*})$ are often called random effects. Besides the individual network structures, researchers aim to estimate the structure and parameters of these fixed effects because they tell us something about the average intraindividual effects.

The random effects can be modeled by assuming a second level normal distribution on all parameters. If the multivariate normal distribution is assumed for all parameters then it is also assumed for the marginal distribution of the means:
\begin{equation}
    \begin{split}
        \bm{\mu}_N &\sim \mathcal{N}(\matr{0}, \bm{\Omega}) \\
        \bm{K}^{\bm{\Omega}} &= \bm{\Omega}^{-1}
    \end{split}
\end{equation}
where $\bm{K}^{\Omega}$ is termed the \textit{between-subjects network}, which is a network between stationary means of different subjects. Thus, estimating the GVAR model on $n > 1$ time series allows the separation of variance into three distinct network structures: (a) temporal networks, (b) contemporaneous networks, and (c) the between-subjects network.

\subsubsection{Pooled and individual LASSO estimation}
First, joint multivariate LASSO estimation with EBIC model selection \cite{abegaz2013sparse} of within-subjects centered data to obtain fixed effects temporal and contemporaneous networks---this I replace with a SUR procedure using FGLS estimation. The variable selection component can be achieved through an initial bootstrap or multi-sample split based procedure which identifies the most stable and important parameters to be included in each network. Second, the GLASSO algorithm with EBIC model selection \cite{foygel2010extended} on sample means of subjects is used to obtain the between-subject network. The first step is repeated for each subject in order to obtain subject-specific networks. The advantages of this approach are that it can estimate fixed effects quickly; scales up to a large number of nodes; can perform model selection for individual networks, and temporal and contemporaneous networks are obtained in the same analysis. The disadvantages are that fixed effects are estimated on pooled data, subject-specific networks are estimated without borrowing information from the multilevel structure, and between-subjects network estimated in a different model. \cite{petersen1983ergodicity,molenaar2009new}.

In the \code{graphicalVAR} package, each variable is first mean-centered (and standardized) across individuals---that is, a grand mean is calculated for each variable and then subtracted from each measurement of that variable; the standardization is simply the default setting in the \code{mlGraphicalVAR} function. Next, subject-specific means are calculated for each individual on each variable. Then, the variables are person-mean centered---this means that the variables are centered \textit{within-subjects}; e.g., the mean of V1 for person 1 is subtracted from each observation of V1 for that person, and so on for all variables, and for all people. Finally, the lagged predictor and response variables are created---again, separately for each person. 

\subsubsection{Two-step frequentist multilevel model}
First, sequential univariate multilevel regression models, with within-subject centered lagged variables as within-subjects level predictors and sample means of all other variables as between-subject predictors. Second, sequential multilevel regression models using the residuals from the first step; residuals of one variable are predicted by residuals of all other variables in the same measurement occasion. Advantages are that the individual models borrow information from the multilevel structure; it scales well up to 8 nodes (with correlated random effects), or 20 nodes (with orthogonal random effects); many random effects variances correlations can be estimated; fast estimation of individual networks. Cons are slow estimation of larger data sets, no model selection, and does not scale well past 20 nodes. 

This method has been extended from the original proposal by \cite{bringmann2013network}. In the first step, variables are within-subject centered, and subject sample means are included as between-subjects predictors. This allows us to estimate the between-subjects network by collecting the univariate regression coefficients and symmetrizing the resulting matrix. In the second step, we predict the residuals from each autoregressive equation using the residuals from all other equations at the same measurement point. Again, these are collected and symmetrized to obtain contemporaneous networks. This is the procedure for \code{mlVAR}:
\begin{enumerate}
    \item Center and standardize all variables
    \item Save person-means for each variable---repeat each for each observation
    \item Lag the variables to create outcomes
    \item Lag the variables to create predictors
    \item Within-person center the predictors
\end{enumerate}
For $N$ individuals who have each been observed on $k$ variables $T$ times, the final response matrix should be $\lbrack N \cdot (T - 1) \rbrack \times k$, and the predictor matrix $\lbrack N \cdot (T - 1) \rbrack \times 2k$.

Then, for the temporal regressions, all lagged predictors are used as random effects and as fixed effects, and then person-means are used as fixed effects, except for whichever variable is the outcome. So, all lagged predictors are always used as fixed and random effects, but then for the person-mean fixed effects we leave out the variable that is being used as the outcome. REML is also not used. This is for \textit{correlated} random temporal effects. Orthogonal random temporal effects means that the random intercept, and each of the random predictor slopes (without intercepts), are specified separately in \code{lmer}. `Fixed' then simply means that there is only a random intercept.

Next, for the contemporaneous networks, we take the residuals from each of the temporal models. We then predict the residuals from one model using the residuals from each other model as both fixed and random effects. There is no fixed nor random intercept in these models, however; only fixed and random effects of the other model residuals. 

The \code{graphicalVAR} package prepares data in a similar way, but with a slight difference:
\begin{enumerate}
    \item Center and standardize all variables
    \item Save person-means for each variable
    \item Within-person center the variables
    \item Lag variables to create outcomes/predictors
\end{enumerate}
The only difference here, really, is that the within-person centering occurs prior to creating lagged variants of the data, rather than centering after lagging. However, this may have been on purpose... with \code{graphicalVAR}, we do not use the within-person means in any of the models except for the between-subjects network. But those means are part of the models in the \code{mlVAR} framework. Hmm...

\subsubsection{N = 1 Graphical VAR}
The lag-1 VAR model is denoted as a regression model on the previous measurement occasion:
\begin{equation}
    \begin{split}
        \matr{y}_t &= \matr{y}_{t-1} \bm{\beta} + \bm{\varepsilon}_t \\
        \bm{\varepsilon}_T &\sim \mathcal{N}(\matr{0},\bm{\Theta})
    \end{split}
\end{equation}
but to indicate that we're explicitly modeling the contemporaneous effects as well, we can write the GVAR model as:
\begin{equation}
    \matr{y}_T\,|\,\matr{y}_{T-1} = \matr{y}_{t-1} \sim \mathcal{N}(\matr{y}_{t-1} \bm{\beta}, \bm{\Theta})
\end{equation}

The GVAR implies the following expression for the variance-covariance matrix of $by_t$ (dropping matrix indexing subscripts for notational clarity):
\begin{equation}
    \begin{split}
        \Var(\by_T) &= \Var(\Beta\by_{T-1} + \beps_T) \\
        \bSigma &= \Beta\bSigma\Beta\tran + \bTheta
    \end{split}
\end{equation}
in which we make the assumption of stationarity and the assumption that residuals $\beps_T$ are uncorrelated with $\by_{T-1}$. Now, we can make use of the vectorization operator Vec and the Kronecker product $\otimes$ to obtain $(\I - \Beta\otimes\Beta)^{-1} \vec{\bTheta} = \vec{\bSigma}$, which gives us an expression of the elements of $\bSigma$ in terms of $\Beta$ and $\bTheta$.
 
\section{Two-Step Multilevel VAR}
We use Roman letters to represent observed variables, and Greek letters to indicate parameters or latent variables. Nonboldface letters indicate a single value---uppercase nonboldface letters indicate a random variable, and lowercase nonboldface letters indicate a realization. $t$ denotes a measurement occasion, and $T$ denotes a random measurement occasion; $i\ (i \in \{1, 2, \ldots, k\})$ to denote a measured variable, and $p\ (p \in \{1,2,\ldots,n\})$ to denote a subject and $P$ to denote a random subject. We use lowercase boldface letters to denote column vectors and uppercase boldface letters to denote matrices. Subscripts will indicate if these are random or fixed. For example, $\bm{B}_p$ will denote a fixed matrix for subject $p$, and $\bm{B}_P$ will denote the matrix of some random subject $P$ (which has a distribution). 

Because we are interested in finding the dynamics among items, we use vector $\matr{y}$ to denote the set of all items. For the observed variables, we use consistent subscripts (measurement, subject) to denote which items are contained in the vector. For instance, $\matr{y}_{t,\,p}$ denotes all responses of subject $p$ at time $t$, and $\matr{y}_{T,\,p}$ denotes all responses of subject $p$ at a random point $T$. A set in this subscript indicates multiple responses. For example, we use $\matr{y}_{\{t - 1,\,t\},\,p}\tran = \lbrack \matr{y}_{t-1,\,p}\tran,\ \matr{y}_{t,\,p}\tran \rbrack$ to denote a set of lagged and current responses from subject $p$ around time point $t$. Then they talk about some $C$ versus $c$ bullshit but it really doesn't make any goddamn sense. Lastly, other subscripts denote subsets of a vector or matrix, with notation $-(\ldots)$ indicating the subset of everything except $\{\ldots\}$.

The primary model to estimate is:
\begin{equation}
    \matr{y}_{T,\,p}\,|\,\matr{y}_{T-1,\,p} = \matr{y}_{t-1,\,p} \sim \mathcal{N} \big( \bm{\mu}_p + \bm{B}_p (\matr{y}_{t-1,\,p} - \bm{\mu}_p), \bm{\Theta}_p \big).
\end{equation}
Specifically, we are interested in estimating the between-subjects network $\bm{K}^{\bm{\Omega}} = \bm{\Omega}^{-1} = \text{Var}(\bm{\mu}_P)^{-1}$ and the (distributions of) temporal networks $\bm{B}_p$ and contemporaneous networks $\bm{K}_p^{\bm{\Theta}} = \bm{\Theta}_p^{-1}$. 

Because the joint conditional distribution of the above model is assumed to be normal, it follows that the marginal distribution of every variable is univariate normal and can be obtained by dropping all other parameters from the distribution
\begin{equation}
    y_{T,\,p,\,i}\,|\,\matr{y}_{T - 1,\,p} = \matr{y}_{t-1,\,p} \sim \mathcal{N} \big( \mu_{p,\,i} + \bm{\beta}_{p,\,i} (\matr{y}_{t-1,\,p} - \bm{\mu}_p ), \theta_{p,\,i} \big),
\end{equation}
in which $y_{T,\,p,\,i}$ denotes the $i$th element of $\matr{y}_{T,\,p}$; $\bm{\beta}_{p,\,i}$ indicates the row vector of the $i$th row of $\bm{B}_p$, and $\theta_{p,\,i}$ denotes the $i$th element of $\bm{\Theta}_p$. When drawn as a temporal network, the edges point to node $i$. Many software packages do not allow the estimation of $\bm{\mu}_p$ as described above. In this case, the sample means of every subject $\bar{\matr{y}}_p$ can be taken as a substitute for $\bm{\mu}_p$ \cite{hamaker2015center}. The model then becomes a univariate regression model with within-subject centered predictors, estimable by functions such as \code{lmer} in \code{lme4} [CITE]. The Level 1 model then becomes:
\begin{equation}
    \begin{split}
        y_{t,\,p,\,i} &= \mu_{p,\,i} + \bm{\beta}_{p,\,i} (\matr{y}_{t-1,\,p} - \bar{\matr{y}}_p) + \varepsilon_{t,\,p,\,i} \\
        \varepsilon_{T,\,p,\,i} &\sim \mathcal{N} (0, \theta_{p,\,i}), 
    \end{split}
\end{equation}
and the Level 2 model becomes:
\begin{equation}
    \begin{bmatrix}
        \mu_{P,\,i} \\
        \bm{\beta}_{P,\,i}
    \end{bmatrix}
    \sim \mathcal{N} \left(
    \begin{bmatrix}
        0 \\
        \bm{\beta}_{*i}
    \end{bmatrix} ,
    \begin{bmatrix}
        \omega_{\mu_{i}} & \bm{\omega}^{(\bm{\beta}_i \mu_i)\tran} \\
        \bm{\omega}^{(\bm{\beta}_i\mu_i)} & \bm{\Omega}^{(\bm{\beta}_i)}
    \end{bmatrix} \right).
\end{equation} 
Estimation of these univariate models requires integrating over a simpler integral than estimation of multivariate models. Thus, sequential estimation using univariate models have been used in estimating multilevel VAR models \cite{bringmann2013network}. A downside is that not all parameters are included in the model---correlations between means (between-subject effects) and between contemporaneous covariances are not retained, as well as the correlations between temporal edges pointing to different nodes. Another downside is that estimating correlated random effects does not work well for models with many predictors. In cases with larger networks (e.g., 20 nodes), we can estimate uncorrelated random effects, termed \textit{orthogonal estimation}. 

The methodology of \cite{bringmann2013network} does not estimate contemporaneous or between-subjects networks. With the two-step method, Step 1 follows the procedure \cite{bringmann2013network} with the addition of between-subjects effects \cite{hamaker2015center}.

\subsection{Step 1: Temporal and Between-Subjects Networks}
To obtain estimates of between-subjects effects, the sample means of every subject $\bmy_p$ can be included as predictors at the subject level (except for the mean of the dependent variable. With this extension, the Level 2 model for the person-specific mean of the $i$th variable now becomes:
\begin{equation}
    \mu_{p,\,i} = \bbeta_i^\mu \bmy_{p,\,-i} + \eps_{p,\,i}^\mu
\end{equation}
in which we use $\bbeta_i^\mu$ to denote the $i$th row (without the diagonal element $i$) of an $m \times m$ matrix $\Beta^\mu$, and $\bmy_{p,\,-i}$ denotes the vector $\bmy_p$ without the $i$th element. Because $\bar{y}_{p,\,i}$ is itself an estimate of $\mu_{p,\,i}$, the Level 2 equation takes the form of a multiple regression model. Thus, these estimates can be used to estimate a GGM between the means \cite{lauritzen1996graphical,meinshausen2006high}---the between-subjects network:
\begin{equation}
    \Kappa^{\bmu} \approx \bDelta^{\bmu} (\I - \Beta^{\bmu}),
\end{equation}
where $\delta_{ii}^\mu = 1/\Var(\eps_{P,\,i})^\mu$. Due to the estimation in a multilevel framework, the resulting matrix will not be perfectly symmetric and must be made symmetric by averaging lower and upper triangular elements (or is this the square-root conversion?). Thus, each edge (i.e., partial correlation) in the between-subjects network is estimated by standardizing and averaging two regression parameters. Obtaining the between-subjects effects using regression coefficients rather than correlating the estimated means leads to standard errors that can be used to select significant edges. 
\subsection{Step 2: Contemporaneous Networks}
The contemporaneous network can be estimated through the residuals of the multilevel model that estimate the temporal and between-subject effects. These residuals can be used to run multilevel models that predict the residuals of one variable from the residuals of other variables at the same time point. Let $\hat{\eps}_{t,\,p,\,i}$ denote the estimated residual variable $i$ at time point $t$ of person $p$, and let $\hat{\eps}_{t,\,p,\,-i}$ denote the vector of residuals of all other variables at this time point. The Level 1 model then becomes:
\begin{equation}
    \hat{\eps}_{t,\,p,\,i} = \bbeta_{p,\,i}^{\bTheta} \hat{\beps}_{t,\,p,\,-i} + \eps_{t,\,p,\,i}^{\bTheta}
\end{equation}
in which $\bbeta_{p,\,i}^{\bTheta}$ represents the $i$th row (without the diagonal element $i$) of an $m \times m$ matrix, $\Beta_p^{\bTheta}$, and $\eps_{t,\,p,\,i}^{\bTheta}$ represents a residual. In the Level 2 model, we again assign a multivariate normal distribution to parameters in $\bbeta_i^{\bTheta}$. The Level 1 equation also takes the form of a multiple regression model. Thus, the process can again be seen as the nodewise GGM procedure:
\begin{equation}
    \Kappa_p^{\bTheta} \approx \bDelta_p^{\bTheta} (\I - \Beta_p^{\bTheta}),
\end{equation}
with $\delta_{p,\,i}^{\bTheta} = 1/\Var(\eps_{T,\,p,\,i}^{\bTheta})$. Again, the matrices can be made symmetric by averaging the upper and lower triangle elements. By using univariate multilevel regressions, rather than simply correlating the residuals, we can impose multilevel structure on the partial correlations in order to estimate fixed and random effects. To threshold the contemporaneous and/or between-subjects networks, we can apply the "AND" or the "OR" rule \cite{barber2015high}.

